%═══════════════════════════════════════════════════════════════════════════════
\chapter{Demostración de la plantilla Catppuccin}
\label{ch:demo}
%═══════════════════════════════════════════════════════════════════════════════

Este capítulo ilustra el uso de la plantilla \texttt{latex-catppuccin-template}.
Todos los elementos mostrados utilizan los colores de la paleta Catppuccin,
adaptándose automáticamente al tema seleccionado (Latte, Frappé, Macchiato o Mocha).

%───────────────────────────────────────────────────────────────────────────────
\section{Entornos teoremáticos}
%───────────────────────────────────────────────────────────────────────────────

La plantilla define varios entornos mediante el generador \texttt{generator.tex}.
Cada uno tiene su propio color de marco y fondo.

\begin{definicion}{Espacio de Hilbert}
Un \textbf{espacio de Hilbert} es un espacio vectorial dotado de un producto
interno que induce una norma completa.
\end{definicion}

\begin{teorema}{Teorema de representación de Riesz}
Sea $\mathcal{H}$ un espacio de Hilbert. Para todo funcional lineal continuo
$\varphi: \mathcal{H} \to \mathbb{R}$ existe un único $y \in \mathcal{H}$ tal que
$\varphi(x) = \langle x, y \rangle$ para todo $x \in \mathcal{H}$.
\end{teorema}

\begin{proposicion}{Identidad de polarización}
En un espacio de Hilbert, el producto interno se puede recuperar de la norma mediante:
\[
\langle x, y \rangle = \frac{1}{4}\left(\|x+y\|^2 - \|x-y\|^2\right).
\]
\end{proposicion}

\begin{ejemplo}{Producto interno en $\mathbb{R}^n$}
El producto interno euclídeo $\langle x, y \rangle = \sum_{i=1}^n x_i y_i$ convierte
a $\mathbb{R}^n$ en un espacio de Hilbert.
\end{ejemplo}

\begin{observacion}{Completitud}
La condición de completitud es esencial: no todo espacio con producto interno es
de Hilbert; se requiere que toda sucesión de Cauchy converja.
\end{observacion}

\begin{fisicaconexion}{Espacios de Hilbert en mecánica cuántica}
En mecánica cuántica, los estados de un sistema físico se representan mediante
vectores en un espacio de Hilbert complejo. La evolución temporal está gobernada
por operadores unitarios.
\end{fisicaconexion}

%───────────────────────────────────────────────────────────────────────────────
\section{Ecuaciones destacadas}
%───────────────────────────────────────────────────────────────────────────────

Puedes escribir ecuaciones con colores semánticos usando los comandos definidos
en \texttt{config/equation-styles.tex} (si los has creado). Aquí un ejemplo con
colores manuales usando \texttt{\textbackslash textcolor}.

\begin{equation}
\label{eq:demo}
\textcolor{CtpBlue}{\int_{\Omega}} \textcolor{CtpPeach}{\nabla u \cdot \nabla v}
\, \textcolor{CtpTeal}{dx} = \textcolor{CtpBlue}{\int_{\partial\Omega}}
\textcolor{CtpPeach}{u \frac{\partial v}{\partial n}} \, \textcolor{CtpTeal}{dS}.
\end{equation}

Y una ecuación más compleja con varios elementos:

\[
\textcolor{CtpMauve}{\frac{\partial \rho}{\partial t}} +
\textcolor{CtpSky}{\nabla \cdot (\rho \mathbf{u})} =
\textcolor{CtpGreen}{0}.
\]

%───────────────────────────────────────────────────────────────────────────────
\section{Algoritmos}
%───────────────────────────────────────────────────────────────────────────────

Los algoritmos se escriben con el paquete \texttt{algorithm2e} y se envuelven
automáticamente en una caja con los colores del tema. Las palabras clave,
comentarios y números de línea están coloreados según la configuración de
\texttt{config/algorithms.tex}.

\begin{algorithm}[H]
\caption{Método de descenso de gradiente}
\label{alg:gradient}
\LinesNumbered
\KwIn{Función $f: \mathbb{R}^n \to \mathbb{R}$, gradiente $\nabla f$,
      punto inicial $x_0$, tasa de aprendizaje $\eta$, tolerancia $\varepsilon$}
\KwOut{Punto aproximado de mínimo $x^*$}
$x \leftarrow x_0$\;
\While{$\|\nabla f(x)\| > \varepsilon$}{
  $x \leftarrow x - \eta \nabla f(x)$\;
}
\Return $x$\;
\end{algorithm}

\begin{algorithm}[H]
\caption{Adam optimizador (paso básico)}
\label{alg:adam}
\LinesNumbered
\KwIn{Gradiente $g$, parámetros $m$, $v$, $t$, $\beta_1$, $\beta_2$, $\eta$, $\varepsilon$}
\KwOut{Actualización $\Delta\theta$, nuevos $m$, $v$, $t$}
$m \leftarrow \beta_1 m + (1-\beta_1) g$\;
$v \leftarrow \beta_2 v + (1-\beta_2) g^2$\;
$\hat{m} \leftarrow m / (1-\beta_1^t)$\;
$\hat{v} \leftarrow v / (1-\beta_2^t)$\;
$\Delta\theta \leftarrow -\eta \hat{m} / (\sqrt{\hat{v}} + \varepsilon)$\;
$t \leftarrow t+1$\;
\Return $\Delta\theta$, $m$, $v$, $t$\;
\end{algorithm}

%───────────────────────────────────────────────────────────────────────────────
\section{Código fuente}
%───────────────────────────────────────────────────────────────────────────────

La plantilla incluye configuración para \texttt{listings} con la paleta Catppuccin.
Aquí hay ejemplos en Python y C++.

\begin{lstlisting}[language=Python, caption={Ejemplo en Python}]
import numpy as np

def relu(x):
    return np.maximum(0, x)

# Red neuronal de una capa
class Perceptron:
    def __init__(self, input_dim, output_dim):
        self.W = np.random.randn(input_dim, output_dim) * 0.01
        self.b = np.zeros(output_dim)

    def forward(self, x):
        return relu(x @ self.W + self.b)
\end{lstlisting}

\begin{lstlisting}[language=C++, caption={Ejemplo en C++}]
#include <iostream>
#include <vector>

template<typename T>
class Matrix {
private:
    std::vector<std::vector<T>> data;
public:
    Matrix(size_t rows, size_t cols) : data(rows, std::vector<T>(cols)) {}
    
    T& operator()(size_t i, size_t j) { return data[i][j]; }
    
    friend std::ostream& operator<<(std::ostream& os, const Matrix& m) {
        for (const auto& row : m.data) {
            for (const auto& val : row) os << val << " ";
            os << "\n";
        }
        return os;
    }
};
\end{lstlisting}

%───────────────────────────────────────────────────────────────────────────────
\section{Ejercicios}
%───────────────────────────────────────────────────────────────────────────────

La plantilla proporciona entornos para ejercicios teóricos y computacionales.

\begin{ejercicio}
Demuestre que la función de pérdida MSE para regresión lineal conduce a las
ecuaciones normales.
\end{ejercicio}

\begin{ejercicio}
Sea $f(x) = e^{-x^2}$. Calcule la transformada de Fourier de $f$.
\end{ejercicio}

\begin{ejerciciocomp}{Implementación de descenso de gradiente}
Implemente el algoritmo \ref{alg:gradient} para la función
$f(x) = x^2 + 2x + 1$ y verifique que converge al mínimo $x^* = -1$.
\end{ejerciciocomp}

\begin{ejerciciocomp}{Simulación de proceso OU}
Simule la ecuación de Langevin $dx = -\theta x\,dt + \sigma dW$ usando
el método de Euler-Maruyama y compare la distribución estacionaria con la
teórica $\mathcal{N}(0, \sigma^2/(2\theta))$.
\end{ejerciciocomp}

%───────────────────────────────────────────────────────────────────────────────
\section{Estrella Polar (opcional)}
%───────────────────────────────────────────────────────────────────────────────

La sección \textbf{Estrella Polar} es opcional y aplica los conceptos del
capítulo a problemas físicos concretos.

\begin{estrella}{Reactor de fisión como proceso estocástico}
La cinética puntual de un reactor nuclear puede modelarse mediante ecuaciones
diferenciales estocásticas que describen la población de neutrones y precursores.
La varianza del flujo neutrónico está relacionada con la reactividad a través
de la fórmula de Campbell.
\end{estrella}

%───────────────────────────────────────────────────────────────────────────────
\section{Dificultad de las secciones}
%───────────────────────────────────────────────────────────────────────────────

Cada sección puede incluir una indicación de dificultad usando el comando
\texttt{\textbackslash dificultad}. Por ejemplo:

\dificultad{3}{2}

Este comando muestra dos valores de dificultad (matemática y física) en una
caja con borde sutil y estrellas.

%───────────────────────────────────────────────────────────────────────────────
\section{Notas bibliográficas}
%───────────────────────────────────────────────────────────────────────────────

Al final de cada capítulo, puedes incluir notas bibliográficas con el entorno
\texttt{notasbib}. Este entorno produce una caja con fondo gris claro que
contiene referencias y comentarios adicionales.

\begin{notasbib}
Para una introducción completa a los espacios de Hilbert, véase
\cite{reed1980functional}. El teorema de representación de Riesz es fundamental
en análisis funcional y ecuaciones en derivadas parciales.
\end{notasbib}